\documentclass[10pt]{article}
\usepackage[T1]{fontenc}
\usepackage[utf8]{inputenc}
\usepackage[french]{babel}
\usepackage{amsmath,amssymb,mathrsfs,mathpazo}
\usepackage{geometry}
\geometry{a4paper,margin=2cm}
\usepackage{graphicx,array,tabularx,tabularray,colortbl,multirow,multicol,booktabs,enumitem,pifont,ulem,fancybox,diagbox,cancel}
\usepackage{pstricks,pstricks-add,pst-plot,pst-tree,pst-node,pst-text,pst-func,pst-math,pst-bspline,pst-3dplot,pst-eucl,pst-solides3d}
\let\pscalculate\undefined
\usepackage{pst-bezier}
\usepackage{tikz,pgfplots,tkz-tab}
\usetikzlibrary{arrows,arrows.meta,patterns,calc,positioning,trees,decorations.markings,shapes.geometric,shapes.misc,shapes.arrows,intersections,decorations.pathmorphing,decorations.pathreplacing,backgrounds,shadows,babel}
\pgfplotsset{compat=1.18}
\usepgfplotslibrary{fillbetween}
\usepackage[np]{numprint}
\usepackage{listings,makecell,pas-tableur,twoopt}
\let\barre\undefined
\usepackage[pstricks]{tablvar}
\definecolor{darkgreen}{RGB}{0,100,0}
\providecommand{\red}{\color{red}}
\providecommand{\blue}{\color{blue}}
\providecommand{\green}{\color{green}}
\providecommand{\black}{\color{black}}
\providecommand{\white}{\color{white}}
\providecommand{\gray}{\color{gray}}
\providecommand{\cyan}{\color{cyan}}
\providecommand{\magenta}{\color{magenta}}
\providecommand{\yellow}{\color{yellow}}
\providecommand{\violet}{\color{violet}}
\providecommand{\orange}{\color{orange}}
\providecommand{\euro}{\texteuro}
\providecommand{\EUR}{\texteuro}
\providecommand{\ieme}{\textsuperscript{e}}
\providecommand{\degres}{\ensuremath{^\circ}}
\providecommand{\degree}{\ensuremath{^\circ}}
\providecommand{\vv}[1]{\overrightarrow{#1}}
\providecommand{\Coord}[1]{#1}
\providecommand{\decosix}{}
\providecommand{\starredbullet}{\ensuremath{\bullet}}
\pagestyle{empty}
\setlength{\parindent}{0pt}
\usepackage{adjustbox}\DeclareUnicodeCharacter{2236}{:}\DeclareUnicodeCharacter{2190}{\ensuremath{\leftarrow}}
\providecommand{\R}{\mathbb{R}}
\providecommand{\N}{\mathbb{N}}
\providecommand{\Z}{\mathbb{Z}}
\providecommand{\Q}{\mathbb{Q}}
\providecommand{\D}{\mathbb{D}}
\providecommand{\C}{\mathbb{C}}
\providecommand{\e}{{\mathrm{e}}}
\providecommand{\vect}[1]{\overrightarrow{#1}}
\providecommand{\vectt}[1]{\overrightarrow{\mathrm{#1}}}
\providecommand{\vvt}[1]{\overrightarrow{\mathrm{#1}}}
\providecommand{\np}[2][]{{#2~#1}}
\providecommand{\vv}[1]{{\overrightarrow{#1}}}
\renewcommand{\d}{{\mathrm{d}}}
\providecommand{\degres}{{^\circ}}
\providecommand{\degre}{{^\circ}}
\providecommand{\blue}{\color{blue}}
\providecommand{\red}{\color{red}}
\providecommand{\enskip}{\quad}
\providecommand{\llbracket}{\lbrack\!\lbrack}
\providecommand{\rrbracket}{\rbrack\!\rbrack}
\providecommand{\euro}{{\text{€}}}
\providecommand{\no}{{\text{n°}}}
\providecommand{\interval}[2]{{[#1;#2]}}
\providecommand{\intervFO}[2]{{[#1;#2[}}
\providecommand{\dotfill}{\ldots}

\providecommand{\ds}{\displaystyle}
\providecommand{\textding}[1]{#1}
\providecommand{\Oij}{\(\mathrm{O};\vec{\imath},\vec{\jmath}\)}
\providecommand{\Oijk}{\(\mathrm{O};\vec{\imath},\vec{\jmath},\vec{k}\)}
\providecommand{\Ouv}{\(\mathrm{O};\vec{u},\vec{v}\)}
\providecommand{\intervFF}[2]{\ensuremath{[#1;#2]}}
\renewcommand{\euro}{\ifmmode\text{\texteuro}\else\texteuro\fi}
\DeclareUnicodeCharacter{00B0}{\ifmmode^{\circ}\else\textdegree\fi}
\DeclareUnicodeCharacter{20AC}{\ifmmode\text{\texteuro}\else\texteuro\fi}
% Énoncés : Éducation nationale. Transcription et corrigés : APMEP, Denis Vergès et contributeurs. Découpage et classement : CoopMaths, Sébastien Lozano, Éric Elter et contributeurs. Adaptation de la présentation : Maths Explorer.
% CC BY-SA : https://coopmaths.fr/www/about/
% https://coopmaths.fr/alea/static/bac/2026/tex/bac_2026_06_sujet2_metropole_4.tex
\begin{document}
\noindent\textit{APMEP / CoopMaths. CC BY-SA. Présentation adaptée par Maths Explorer.}\par\medskip
\subsection*{\centering Partie A}

			Soit $f$ une fonction définie et dérivable sur l'intervalle $[0 ;+\infty[$ dont on donne ci-dessous une partie de la courbe représentative $\mathcal{C}_f$ dans un repère orthogonal.

			\begin{minipage}{0.45\linewidth}
				On note $f'$ sa fonction dérivée.\\

			\begin{list}{\textbullet}{On précise que:}
				\item la courbe $\mathcal{C}_f$ admet une tangente parallèle à l'axe des abscisses au point C d'abscisse 1;

				\item la tangente $T_{\mathrm{A}}$ à la courbe $\mathcal{C}_f$ au point A$\left(\dfrac{1}{2} ; 0 \right)$ passe par le point B$\left(1 ; \e^{2}\right)$.
			\end{list}

\medskip

			\begin{enumerate}
				\item À l'aide du graphique ci-contre :
					\begin{enumerate}
						\item Donner $f'(1)$.

						\item Donner la solution de l'équation $f(x)=0$ sur l'intervalle $[0 ; 3]$.
					\end{enumerate}

				\item Déterminer $f'\left(\dfrac{1}{2}\right)$.
			\end{enumerate}
			\end{minipage}
			\hfill
			\begin{tikzpicture}[baseline={(current bounding box.center)}]
				\begin{axis}[/pgf/number format/.cd, use comma,
					x={15mm}, y={7.5mm},
					xmin = -0.4, xmax = 4.7,
					ymin = -3.4, ymax = 8.6,
					xtick={0,1,...,5},
					ytick ={-3,...,8},
					xticklabels={0,1},
					yticklabels={,,,0,1},
					minor tick num = 0,
					grid=major, axis lines =center]
					\draw[teal,line width=0.8pt] (0,-7.389)--(1.5,14.778) node[pos=0.71,right]{$T_{\mathrm{A}}$};
					\addplot [line width=1.2pt,color=red,smooth,samples=100,domain= 0.2:5 ]{(2*x-1)*exp(-2*x+3)};
					\node[below left] at (axis cs: 0,0) {{0}};
					\node[red] at(axis cs: 3.2,0.6) {$ \mathcal{C}_f $};
					\draw[shift={(axis cs: 0.5,0)}] (2pt,-2pt)--(-2pt,2pt) (-2pt,-2pt)--(2pt,2pt) node[above left]{A} ;
					\draw[shift={(axis cs: 1,7.389)}] (2pt,-2pt)--(-2pt,2pt) (2pt,2pt)--(-2pt,-2pt) node[left]{B} ;
					\draw[shift={(axis cs: 1,2.718)}] (2pt,-2pt)--(-2pt,2pt) (-2pt,-2pt)--(2pt,2pt) node[above right]{C} ;
				\end{axis}
			\end{tikzpicture}

			\begin{enumerate}[start=3]
				\item Parmi les trois courbes $\mathcal{C}_{1}$, $\mathcal{C}_{2}$ et $\mathcal{C}_{3}$ ci-dessous, deux représentent des primitives de la fonction $f$.

				Indiquer lesquelles et justifier.

				\begin{center}
					\begin{tikzpicture}[baseline={(current bounding box.center)}]
						\begin{axis}[/pgf/number format/.cd, use comma,
							x={19mm}, y={10mm},
							xmin = -0.4, xmax = 4.7,
							ymin = -3.4, ymax = 5.4,
							xtick={0,1,...,5},
							ytick ={-3,...,8},
							xticklabels={0,1},
							yticklabels={,,,0,1},
							minor tick num = 0,
							grid=major, axis lines =center]
							\node[below left] at (axis cs: 0,0) {{0}};
							\addplot [line width=1.2pt,color=red,smooth,samples=100, dashed,domain= 0:5 ]{-x*exp(-2*x+3)+exp(1)+1};
							\node[red] at(axis cs: 2.3,3.6) {$ \mathcal{C}_1 $};
							\addplot [line width=1.2pt,color=teal,smooth,samples=100, dotted,domain= 0:5 ]{-x*exp(-2*x+3)+exp(1)-1};
							\node[teal] at(axis cs: 2.3,1.6) {$ \mathcal{C}_2 $};
							\addplot [line width=1.2pt,color=blue,smooth,samples=100, dashdotted ,domain= 0:5 ]{-(x-0.4)*exp(-2*(x-.4)+3)+exp(1)-1};
							\node[blue] at(axis cs: 2.3,0.4) {$ \mathcal{C}_3 $};
						\end{axis}
					\end{tikzpicture}
				\end{center}
			\end{enumerate}


		\subsection*{\centering Partie B}

			\begin{minipage}{0.55\linewidth}
				Un industriel doit créer un logo représentant la forme ci-contre (Fig. 1).

				La dimension de \np[cm]{0,3} indiquée sur la figure est une contrainte imposée par le client.

				\bigskip

				L'industriel choisit de modéliser le contour de ce logo à l'aide d'une portion de la courbe de la fonction $f$ vue en \textbf{partie A}, et de l'image de cette portion par la symétrie par rapport à l'axe des abscisses, sur un intervalle $[0,5 ; \alpha]$.

				Le nombre $\alpha$ est un réel supérieur à 1 que l'on déterminera dans cette partie.

\medskip

				L'ensemble est représenté graphiquement dans un repère orthonormé (Fig. 2) d'unité 1 cm.
			\end{minipage}
			\hfill
			\begin{tikzpicture}[baseline={(current bounding box.center)},>=stealth]
				\begin{axis}[/pgf/number format/.cd, use comma,
					x={12mm}, y={12mm},
					xmin = 0.4, xmax = 4.2,
					ymin = -3, ymax = 3,
					axis lines=none,
					xtick=\empty, ytick=\empty, grid=none]
					\addplot [line width=1.2pt,color=red,smooth,samples=100,domain= 0.5:3.312]{(2*x-1)*exp(-2*x+3)};
					\addplot [line width=1.2pt,color=red,smooth,samples=100,domain= 0.5:3.312]{-(2*x-1)*exp(-2*x+3)};
					\draw[red,line width=1.2pt](axis cs: 3.312,-0.15)--(axis cs: 3.312,0.15) ;
					\draw[line width=0.8pt,<->](axis cs: 3.45,-0.15)--(axis cs: 3.45,0.15) node[pos=0.5,right] {\np[cm]{0,3}};
					\node at(axis cs: 2.3,-2.5) {Fig. 1};
				\end{axis}
			\end{tikzpicture}

			On admet que la fonction $f$ est définie sur l'intervalle $[0 ;+\infty[$ par :
			\[f(x)=(2 x-1) \e^{-2 x+3}.\]
			\begin{minipage}{0.55\linewidth}
			\begin{enumerate}
				\item \begin{enumerate}
					\item Vérifier que pour tout réel $x$ appartenant à l'intervalle $[0 ;+\infty[$ :
					\[f(x)=\e^{2} \times \dfrac{2 x-1}{\e^{2 x-1}}.\]

					\item En déduire la limite de $f$ en $+\infty$.

					(On pourra poser $X = 2x-1)$.
				\end{enumerate}

				\item \begin{enumerate}
					\item Montrer que,

					 pour tout $x$ de l'intervalle $[0 ; +\infty[$,
					\[f'(x)=(-4 x+4) \e^{-2 x+3}.\]
					\item Construire le tableau complet des variations de la fonction $f$ sur l'intervalle $[0 ;+\infty[$.
				\end{enumerate}
				\item Justifier que la contrainte imposée par le client se traduit par l'égalité
				\quad $f(\alpha)=0,15$.

				Montrer qu'il existe une unique valeur de $\alpha$ dans l'intervalle $[1 ;+\infty[$ qui permette de vérifier cette égalité.

				En donner une valeur arrondie au dixième.
			\end{enumerate}
			\end{minipage}
			\hfill
				\begin{tikzpicture}[baseline={(current bounding box.center)},>=stealth]
					\begin{axis}[/pgf/number format/.cd, use comma,
						x={14mm}, y={14mm},
						xmin = -0.4, xmax = 3.8,
						ymin = -3.1, ymax = 3.1,
						xtick={0,1,...,5},
						ytick ={-3,...,8},
						xticklabels={0,1},
						yticklabels={,,,0,1},
						minor tick num = 0,
						grid=major, axis lines =center]
						\node[below left] at (axis cs: 0,0) {{0}};
						\addplot [name path = Cf,line width=1.2pt,color=red,smooth,samples=100,domain= 0.5:3.312]{(2*x-1)*exp(-2*x+3)};
						\addplot [name path = Cg,line width=1.2pt,color=red,smooth,samples=100,domain= 0.5:3.312,dashed]{-(2*x-1)*exp(-2*x+3)};
						\addplot[fill=gray,fill opacity=0.1] fill between[of=Cf and Cg];
						\node[right] at(axis cs: 3.312,-0.2) {$\alpha$};
						\node[red] at(axis cs: 2.7,0.7) {$ \mathcal{C}_f $};
						\node[fill=white] at(axis cs: 2,-3) {Fig. 2};
					\end{axis}
				\end{tikzpicture}

		\subsection*{\centering Partie C}

			On fabrique un porte-clé en découpant une plaque d'acier suivant la forme du logo représenté par la surface grisée en figure 2 de la \textbf{partie B}.

			On précise que cette plaque a une épaisseur de $\np[cm]{0,2}$.

			\begin{enumerate}
				\item On considère :
				\[I=\int_{0,5}^{3,3} f(x) \,\mathrm{d} x\]

				À l'aide d'une intégration par parties, montrer que l'intégrale $I$ arrondie au dixième est égale à 3,6.

				\item On précise que le volume du porte-clé est égal à :

\begin{center}
\emph{Aire du logo} $\times$ \emph{épaisseur de la plaque}
\end{center}

				Déterminer le volume du porte-clé en $\mathrm{cm}^{3}$, arrondi au dixième.
			\end{enumerate}

\end{document}
